\documentclass[11pt]{article}

\usepackage[margin=1in]{geometry}
\usepackage{amsmath}
\usepackage{array}
\usepackage{booktabs}
\usepackage[colorlinks=true,allcolors=blue]{hyperref}
\usepackage{longtable}
\usepackage{microtype}
\usepackage[numbers,sort&compress]{natbib}
\usepackage[T1]{fontenc}

\newcolumntype{P}[1]{>{\raggedright\arraybackslash}p{#1}}

\hypersetup{
  pdftitle={Climate preconditioning and the mechanics of hazardous alpine mass movements: a catalog protocol},
  pdfauthor={Bradley P. Lipovsky}
}

\title{Climate preconditioning and the mechanics of hazardous alpine mass movements:\\a catalog and attribution protocol}
\author{Bradley P. Lipovsky}
\date{Working manuscript, 30 August 2026}

\begin{document}
\maketitle

\begin{abstract}
Large rock, ice, and mixed-material collapses can transform into far-traveled
debris flows or locally extreme tsunamis. Warming changes glacier geometry,
permafrost temperature, meltwater supply, and precipitation phase. However,
event reports seldom separate these preconditioning changes from immediate
weather, earthquakes, structure, and progressive damage. Here, I define a
climate-blind event inventory and an attribution protocol that preserves this
distinction. Two structured frames contain 40 and 38 rows that cross an
operational volume, travel, or runup prefilter. An exact provenance crosswalk
gives every row a visible disposition, and supplemental searches extend the
screen to 148 candidate occurrences. Two independent screens include 73,
retain 68 as uncertain, and exclude 7; 53 included failures have day-scale or
better timing for a later trigger analysis. Trigger, cascade, and repeat-site
dependence remain separate. I then specify detrended seasonal controls,
matched non-failing slopes, mechanism-specific glacier and permafrost
observations, and multiplicity control. Uneven observation and unresolved
primary sources prevent these counts from estimating a frequency trend. No
systematic climatic association is estimated here.
\end{abstract}

\section{Introduction}

Alpine mass movements become especially hazardous when an initial failure
changes material or enters a confined channel. The 2021 Chamoli rock and ice
avalanche transformed into a debris flow and flood that scoured valley walls
and struck two hydropower projects \citep{Shugar2021}. Fjord failures instead
transfer momentum to water. The September 2023 Dickson Fjord rockslide
generated near-source runup of about 200~m and a very-long-period signal that
persisted for nine days \citep{Svennevig2024,Monahan2025}. The August 2025
Tracy Arm landslide produced a measured maximum runup of
$481\pm7.6$~m \citep{Shugar2026}. A useful catalog must describe the initial
failure separately from the downstream cascade.

The mechanical origin of the initial failure remains the central uncertainty.
Warming may remove glacier support from a rock slope, thaw ice in fractures,
or change meltwater pressure beneath a glacier. Weather may provide a shorter
perturbation through rain, snow loading, melt, or rapid temperature change.
Lithology, adverse structure, progressive damage, valley geometry, glacier
dynamics, and earthquakes remain competing or interacting explanations. A
climatic trend can therefore precondition a failure without being its
immediate trigger, and a warm event year does not establish causation.

Here, I use a literature seed to define a reproducible catalog and the tests
that would be needed to evaluate systematic association with climate. The
seed is a protocol-development set, not an inference sample. A separate,
climate-blind discovery inventory supplies the candidate sampling frame. The
two are kept separate so that examples chosen because they motivated a climate
hypothesis cannot enter an association test without being rediscovered under
the same rules as every other event.

\section{Catalog construction and sampling protocol}

The catalog consists of three linked tables. The event table records identity,
group, analysis role, date, location, setting, initial failure, and downstream
transformations. Dates and coordinates have field-specific sources. The
measurement table stores one quantity per row and distinguishes a point value,
approximation, lower bound, upper limit, or range. It records uncertainty type,
observation period, scenario, evidence kind, and source locator separately.
The claim table labels each statement as observation, preliminary observation,
reconstruction, model output, published inference, or project interpretation.
Each event has explicit process-chain, immediate-trigger, and
climate-preconditioning claims, including an unresolved claim when evidence is
absent.

The registered geophysical inclusion rule requires rapid failure of rock, glacier
ice, or a mixed mass in an alpine or fjord setting, a documented time and
source location, and at least one of three operational thresholds: initial
volume of $10^6$~m$^3$, source-to-deposit travel of 2~km, or landslide-tsunami
field runup of 10~m. Downstream damage and population exposure are stored as
consequences but do not determine geophysical inclusion. Snow avalanches
without glacier-ice or rock failure and floods without an initial mass
movement are outside the scope. These thresholds were frozen before the
climate-blind screen. A separately labeled active
slope may enter the protocol-development set when measured deformation and
published failure geometries exist. Barry Arm is such a prospective case, not
an occurred event \citep{Dai2020,Schaefer2023,Barnhart2021,Collins2024}.

Discovery covers 1 January 2000 through 30 August 2026. Two enumerated frames
are the 354-case PANGAEA landslide-tsunami catalog
\citep{Dohmen2025Data} and the 60-event High Mountain Asia rock- and
ice-avalanche inventory \citep{Zhong2024,Zhong2024Data}. A 51-event global
review supplies a cross-check \citep{Mani2023}; its linked spreadsheet was not
enumerated because automated retrieval returned HTTP 403. Independent query
suites for fjord and lake tsunamis, glacier failures, and rock or mixed
avalanches supplement dates and regions not covered by the structured frames.
All exact queries, execution times, and resource classes are archived with the
catalog. Searches and candidate fields omit asserted climate cause and event
climate exposure. Inventory cause columns are ignored.

\begin{table}[t]
\centering
\caption{Cause-blind flow through the two enumerated discovery frames.
``Retained'' applies the occurrence window and a usable event record; the
numeric prefilter uses only published volume, travel, or runup fields. It does
not establish material or setting eligibility.}\label{tab:frames}
\begin{tabular}{@{}lrrrrr@{}}
\toprule
Frame & Published & Downloaded & Window & Retained & Numeric \\
\midrule
PANGAEA landslide tsunamis & 354 & 355 & 87 & 84 & 40 \\
High Mountain Asia & 60 & 60 & 40 & 40 & 38 \\
\bottomrule
\end{tabular}
\end{table}

Every candidate promoted to detailed screening has two differently named
reviewers. Each reviewer sees identity, time, place, initial material and
process, the threshold observation, and its source. Allowed decisions are
include, exclude, or uncertain. An inventory value can discover an event, but
final inclusion requires a primary event reconstruction or authoritative
agency record. A secondary-only record remains uncertain. Conflicting dates,
coordinates, volumes, and aliases remain in the notes rather than being
silently reconciled. Detection era, sensor coverage, and reporting source
must remain in the later observation model because discovery is not stationary.

The reconciliation crosswalk links all 78 numeric-prefilter rows in
Table~\ref{tab:frames} to candidate records. The source bytes, checksums, and
qualifying-row manifest are archived with the catalog. Compound inventory rows
are
expanded into individual occurrences. A row-level size estimate is not
assigned to each component unless the primary source supports that allocation;
otherwise each component remains uncertain. Independent searches identified,
for example, the November 2025 Bartuytsete ice-rock avalanche, whose published
3.5~km runout passes the travel threshold \citep{Dokukin2026}. The resulting
table contains 148 candidates: 73 included, 68 uncertain, and 7 excluded.
Fifty-three included failures have day-scale or better timing and enter the
future trigger-time risk set; the remainder can inform only site-level or
preconditioning analyses.

The seed records were selected because they motivated the study and still
cannot support a frequency trend, an event-conditioned climate distribution,
or a global susceptibility model. The discovery table is closer to an
inference frame, but it also cannot support those estimates until the
registered search cutoff, both screens, disagreement resolution, and catalog
commit are frozen. No atmospheric anomaly or cryospheric exposure is extracted
before that freeze.

The protocol has two estimands. The episode-level trigger estimand compares
conditions at failure with time controls; the site-level preconditioning
estimand compares source slopes with an incidence-density risk set. The two
Aru glaciers are separate episode candidates but share a site-level cluster.
Distinct Aysen, Denali, and Kunlun source failures remain individual events
even though each set shares an earthquake trigger. The unresolved October 2023
Dickson Fjord episode is linked to the September site but receives no analysis
role until its source-specific threshold is established. Barry Arm is excluded
from event-time inference because catastrophic failure has not occurred; it
instead supplies a prospective case for testing susceptibility variables.

\section{Mechanics represented by the seed records}

The two Aru glacier tongues detached after surge-like acceleration on
low-angle terrain; the published mean Aru-1 surface slope is approximately
$13^\circ$. Remote sensing and thermomechanical modeling indicate
that reduced basal friction shifted resistance toward frozen tongues and
lateral margins until those forces were exceeded
\citep{Kaab2018,Gilbert2018}. Gilbert and colleagues inferred a five- to
six-year increase in melt and rain rather than exceptional seasonal forcing.
For Aru-1, approximately 10--25~mm of precipitation during the two days before
failure is a possible short perturbation; Aru-2 has no resolved immediate
trigger. The two collapses traveled approximately 6 and 5~km beyond their
respective glacier termini \citep{Kaab2018}.

At Chamoli, $26.9\times10^6$~m$^3$ of rock and glacier ice failed
from Ronti Peak and transformed into a highly mobile debris flow
\citep{Shugar2021}. The reported 95\% reconstruction interval is
$26.5$--$27.3\times10^6$~m$^3$. The mass was about 80\% rock and 20\% ice;
these are model-dependent composition estimates rather than direct material
measurements. Mapping and reconstruction do not require a glacial-lake
outburst to explain the flood cascade. Geologic structure and progressive
damage remain alternatives to thermal and hydrologic preconditioning. Snow
and ice loading on the headwall and a preceding positive temperature anomaly
have been proposed as final perturbations, but a unique immediate trigger was
not isolated \citep{VanWykDeVries2022}.

At Dickson Fjord, a foliation-parallel rockslide struck a gully glacier before
entering the fjord. The reconstructed rock volume was
$25.5\pm1.7\times10^6$~m$^3$, with a separate
$2.3\pm0.15\times10^6$~m$^3$ entrained-ice estimate
\citep{Svennevig2024}. Rockslide, glacier impact, tsunami, and the
very-long-period signal are observations or reconstructions. The cross-fjord
seiche was first inferred by wave and seismic modeling; sparse SWOT altimetry
later provided direct water-level observations consistent with its phase and
period \citep{Monahan2025}. Multidecadal glacier lowering supports a
debuttressing hypothesis, but no proximate trigger or formal anthropogenic
event attribution was demonstrated.

Tracy Arm and Barry Arm illustrate occurred and prospective consequences of
glacier retreat along fjords. The mapped Tracy Arm volume gives a lower limit
of $63.5\times10^6$~m$^3$ with a reported uncertainty of
$2.7\times10^6$~m$^3$; unmeasured submarine failure may increase the total.
Several days of accelerating microseismicity preceded failure
\citep{Shugar2026}. Regional summer warming since 1875 was
quantitatively attributed to anthropogenic forcing. Observed retreat of South
Sawyer Glacier supports debuttressing as a mechanical inference, but the
glacier retreat itself was not assigned an anthropogenic counterfactual
fraction. At Barry Arm, slope displacement and glacier retreat were temporally
coincident during 2010--2017 \citep{Dai2020}. Their reported correlation of
0.99 compares cumulative trending series and is not by itself a causal test.
Published geometry scenarios span $290$--$689\times10^6$~m$^3$
\citep{Barnhart2021}, while a
four-layer seismic inversion estimates a local mean failure-surface depth of
$188\pm9$~m \citep{Collins2024}. These are uncertain slope geometries, not a
failed volume.

\begin{longtable}{@{}P{0.15\textwidth}P{0.16\textwidth}P{0.12\textwidth}P{0.23\textwidth}P{0.22\textwidth}@{}}
\caption{Seven focal records in the protocol seed. The dependent October 2023
Dickson Fjord episode remains in the machine-readable catalog but is omitted
here. Volumes describe an initial failed mass except at Barry Arm, where the
range describes possible unstable-slope geometries. Shown uncertainties are
reproduced from the sources; their distinct confidence, conservative-sum,
or unspecified semantics are stored in the machine-readable catalog.}\label{tab:events}\\
\toprule
Event and date & Initial failure & Volume (10$^6$ m$^3$) & Cascade & Climate evidence \\
\midrule
\endfirsthead
\toprule
Event and date & Initial failure & Volume (10$^6$ m$^3$) & Cascade & Climate evidence \\
\midrule
\endhead
Nepal, 26 Aug 2026 & Provisional glacial-collapse classification & unresolved & Debris flow and flood traveled about 100 km & No event-specific attribution; process reconstruction remains provisional \citep{USGS2026Nepal,USGS2026NepalSeismic} \\
Chamoli, 7 Feb 2021 & Rock and glacier ice & 26.9 [26.5, 27.3] & Avalanche to debris flow and flood & Several thermal effects are plausible; no single trigger isolated \citep{Shugar2021,VanWykDeVries2022} \\
Aru-1, 17 Jul 2016 & Glacier detachment & $68\pm2$ & Ice avalanche; 6 km beyond terminus & Multiyear melt and rain inferred to increase subglacial water pressure \citep{Gilbert2018} \\
Aru-2, 21 Sep local time 2016 & Glacier detachment & $83\pm2$ & Two pulses; 5 km beyond terminus & Shared multiyear preconditioning; immediate trigger unresolved \citep{Kaab2018,Gilbert2018} \\
Dickson Fjord, 16 Sep 2023 & Foliation-parallel rockslide & $25.5\pm1.7$ rock & Glacier entrainment, tsunami, and nine-day signal & Glacier thinning inferred to debuttress the toe; proximate trigger unresolved \citep{Svennevig2024} \\
Tracy Arm, 10 Aug 2025 & Bedrock landslide & $\geq63.5$ mapped & Glacier/fjord impact, tsunami, and signal lasting at most 36 hours & Regional warming attributed; retreat--debuttressing link is a mechanical inference \citep{Shugar2026} \\
Barry Arm, active & Slow bedrock landslide & 290--689 possible & Potential rapid failure and tsunami & Retreat and motion are temporally coincident; alternatives require testing \citep{Dai2020,Barnhart2021} \\
\bottomrule
\end{longtable}

\section{Mechanical predictors beyond reanalysis}

Atmospheric fields do not measure glacier support, permafrost fracture state,
or basal effective pressure. The analysis will therefore pair reanalysis with
annual optical and radar mapping of terminus position, glacier-slope contact,
surface elevation, and source-slope deformation. Define $\Delta h_i$ as index-
date thickness minus baseline thickness at the slope contact. A signed
stress-scale proxy for support loss is
\begin{equation}
  S_i = -\rho_i g\Delta h_i,
\end{equation}
where $\rho_i$ is ice density, so thinning gives positive $S_i$ and thickening
gives negative $S_i$. This proxy has stress units but is not the lateral
buttressing traction. Geometry-specific elastic or limit-equilibrium
calculations are required to translate it into a support change. Results will
therefore report signed thickness change, its magnitude, mapped contact loss,
and $S_i$ before interpreting debuttressing.

For a first glacier-detachment screening bound, treat the bed as Coulomb till:
\begin{equation}
\begin{aligned}
  N &= p_i-p_w, & \tau_y &= c + N\tan\phi, \\
  \tau_b &\leq \tau_y, & \tau_d &\simeq \rho_i g h\sin\alpha,
\end{aligned}
\end{equation}
where $N$ is effective pressure, $p_i$ is ice overburden pressure, $p_w$ is
water pressure, $c$ and $\phi$ are effective till cohesion and friction angle,
$h$ is ice thickness, and $\alpha$ is surface slope. In a slab that neglects
lateral and longitudinal resistance, basal yield cannot balance driving stress
when
\begin{equation}
  N < N_{\mathrm{crit}} = \frac{\tau_d-c}{\tan\phi}.
\end{equation}
This is a conditional screening range, not an inferred Aru effective pressure;
the analysis must span plausible $c$, $\phi$, geometry, and neglected margin
stresses. Melt or rain can raise $p_w$ and reduce resistance when drainage is
inefficient. The same input can enlarge efficient drainage, lower $p_w$, and
increase resistance.

Permafrost hypotheses require rock-temperature or fracture-ice observations,
not air temperature alone. Where boreholes are absent, the protocol will use
modeled rock-surface temperature, topographic shading, fracture orientation,
and cumulative thermal history as proxies and will label the inference.
Adverse structure is a baseline confounder; deformation may be a mediator or
competing-process observation. The causal specification will determine which
variables are adjusted rather than treating every measured field as a
covariate.

\section{Event-centered climate test}

For each frozen event candidate, reanalysis will provide near-surface air
temperature, precipitation phase and amount, snow accumulation, radiative and
turbulent fluxes, and surface pressure. Prespecified reductions will include
cumulative positive degree days, liquid-water input, recent snow load,
freeze--thaw crossings, and a multidecadal temperature trend. Windows,
directions, and mechanism-specific primary contrasts will be registered before
event values are examined.

Immediate-trigger anomalies and long-term preconditioning require different
controls. For trigger variables, I will remove the local seasonal cycle and a
smooth long-term trend using data outside the event window. The event residual
will be compared with time-matched, same-season blocks. Block resampling will
preserve serial correlation, and the two Aru records will share a site-level
cluster. This avoids a null that ranks recent event years against an earlier,
cooler climate and then interprets the resulting rank as unusual weather.

For persistent preconditioning, the comparison will use incidence-density
risk sets. At each event date, eligible controls are catalog-search source
slopes that remain under observation and have not yet failed. Controls receive
the case index date, and every exposure window ends before that date; a control
may later become a case. Follow-up ends at failure, loss of adequate sensor
coverage, or the registered study end. Matching uses baseline region,
elevation, aspect, relief, lithology where available, and pre-exposure sensor
coverage. It does not use index-date glacier contact or deformation, which may
lie on the causal pathway. A preregistered causal diagram will distinguish
baseline confounders from exposures and mediators.

A hierarchical model will estimate associations while including region and
site-group effects, detection era, and monitoring intensity. Glacier-
detachment sites will have a signed liquid-water/effective-pressure contrast;
glacier-contact rock slopes will have a signed contact-loss/support contrast;
and permafrost rock slopes will have a cumulative thermal contrast. Predictors
will be standardized against their own risk-set distributions and will not be
imputed across inapplicable mechanism classes. Family-wise error will be
controlled across these primary contrasts; secondary predictors will use a
false-discovery-rate procedure. An association under this design would still
not constitute a complete anthropogenic event attribution.

Global products smooth the elevation and aspect that govern steep-terrain
energy balance. Downscaling will begin with corrections tied to a named bias:
temperature adjusted between grid-cell and source elevation using a documented
lapse rate, precipitation separated by phase, and shortwave radiation
corrected for slope, aspect, and shading. Results will be reported before and
after each correction. Regional modeling is warranted only when these
corrections change an event anomaly enough to affect the prespecified test.

\section{Competing hypotheses and discriminating observations}

The analysis will evaluate signs and timing that distinguish mechanisms.
Debuttressing predicts a deformation change point after glacier thinning or
contact loss, with a geometry-dependent lag. Loading has no general sign. For
an added load $\Delta q$ on a planar slope, the simple components are
$\Delta\tau=\Delta q\sin\alpha$ and
$\Delta\sigma_n=\Delta q\cos\alpha$; the Coulomb tendency
$\Delta\tau-\Delta\sigma_n\tan\phi$ changes sign with slope and friction
angle, before crack pressure and three-dimensional geometry are considered.
Hydrologic weakening predicts velocity or
failure probability increasing with water input when drainage is inefficient;
hydrologic strengthening predicts a negative relation or seasonal hysteresis
as drainage evolves. Permafrost degradation predicts deformation associated
with cumulative rock-temperature history and ice-bearing fractures rather
than a single warm day. Progressive damage predicts accelerating deformation
or microseismicity without requiring unusual weather. Structural release
predicts failure along mapped adverse discontinuities. Seismic forcing
predicts temporal coincidence with an earthquake capable of producing
materially large shaking at the source.

Barry Arm provides a prospective case for checking several predictions. Change points
and lags in slope velocity, glacier-contact loss, precipitation, groundwater,
and seasonal temperature can be tested with autocorrelated errors. If a
glacier term loses explanatory power after detrending or after internal creep
state is included, the cumulative 0.99 correlation should not be treated as
evidence for a dominant causal pathway.

\section{Uncertainty and next steps}

Source uncertainty enters at three levels. First, reported intervals are not
interchangeable. Chamoli gives a 95\% reconstruction interval; Aru reports a
conservative linear sum of DEM and post-event correction errors; the
probability semantics of some reported Dickson and Tracy uncertainties remain
unspecified and are labeled \texttt{reported\_unspecified}. The Barry Arm depth spread describes agreement
among tested parameterizations and time periods, not full inversion
nonuniqueness. Second, source classifications can change. The
reviewed Nepal seismic record now calls the source a glacial collapse and
debris flow, whereas the dated event page retains broader provisional wording
\citep{USGS2026Nepal,USGS2026NepalSeismic}. Both records are preserved. Third,
some quantities are conditional scenarios. Barry Arm volumes and tsunami
heights describe possible geometries, not observations or forecasts of a
failure time.

This version completes the climate-blind frame reconciliation and double
screening. The unresolved records require primary-source follow-up and the
registered terrain-mask check, but they are not silently removed. The next
analysis must freeze this candidate version, acquire matched non-failing
slopes, and calculate glacier-contact histories, deformation records, and
permafrost proxies before inspecting event climate anomalies. Only after those
data and the statistical specification are frozen can the study estimate
whether recent hazardous alpine mass movements occupy systematically different
climatic or cryospheric conditions.

\section{Conclusions}

The catalog establishes a traceable distinction among initial failure,
downstream cascade, climate preconditioning, and immediate trigger. The
climate-blind screen retains 73 eligible occurrences while keeping 68
unresolved cases visible. It also shows why reanalysis alone is insufficient:
glacier detachments require effective-pressure bounds, while rock-slope cases
require glacier geometry, support change, structure, and deformation.
Detrended seasonal controls, matched non-failing slopes, and explicit competing
predictions are the next tests. The present contribution is the sampling frame
and protocol, not evidence that warming systematically triggered these events.

\bibliographystyle{plainnat}
\bibliography{references}

\end{document}
